\section{بار کاری حافظه (\lr{memory})}
\fig[0.55\linewidth]{fig/memory.png}{توان عملیاتی (\lr{throughput}) در آزمون \lr{memory}}
\begin{center}
\begin{tabular}{lrrr}
\toprule
محیط & \lr{Throughput (MiB/s)} & انحراف معیار & \lr{Total time (s)} \\
\midrule
\lr{Bare Metal (WSL)} & \lr{1388.73} & \lr{63.93} & \lr{$\approx$10.00} \\
\lr{VM} & \lr{1385.12} & \lr{93.58} & \lr{$\approx$10.00} \\
\lr{Docker} & \lr{1384.78} & \lr{66.62} & \lr{$\approx$10.00} \\
\bottomrule
\end{tabular}
\end{center}
\textbf{تحلیل.} میانگین سه محیط تا کمتر از ۰٫۳٪ با هم برابرند، در حالی که پراکندگی بین تکرارهای یک محیط (حدود ۵ تا ۱۰٪) بسیار بزرگ‌تر از اختلاف بین محیط‌هاست؛ پس هیچ تفاوت معنادار قابل‌تشخیصی وجود ندارد. زمان کل (\lr{Total time}) نیز در همهٔ اجراها حدود ۱۰ ثانیه است، یعنی آزمایش با محدودیت زمانی پیش‌فرض \lr{sysbench} تمام شده (حجم ۲۰ گیگابایتی آزمون در حافظهٔ \lr{VM} جا می‌شود) و معیار مقایسه همان \lr{throughput} است. در \lr{VM} ترجمهٔ آدرس حافظه با پشتیبانی سخت‌افزاری (\lr{EPT/NPT}) انجام می‌شود و پس از پر شدن \lr{TLB} سربار کمی دارد؛ در \lr{Docker} نیز حافظه مستقیماً از همان هسته مدیریت می‌شود. بنابراین هر دو سربار ناچیزی برای دسترسی به حافظه دارند. (در آزمون تصادفی، سربار \lr{VM} می‌تواند در صورت \lr{TLB miss} زیاد بیشتر شود، که در این داده‌ها دیده نشد.)
